\documentclass{article}
\usepackage{amsmath, amssymb, amsthm}
\usepackage{hyperref}
\usepackage{geometry}
\geometry{margin=1in}

\title{Erd\H{o}s Problem \#450}
\date{Last updated: 2025-08-31}
\author{Source: erdosproblems.com}

\begin{document}
\maketitle

\noindent\textbf{Status:} OPEN \quad \textbf{No prize}

\noindent\textbf{Tags:} number theory, divisors

\medskip

\noindent\textbf{Problem Statement:}

How large must $y=y(\epsilon,n)$ be such that the number of integers in $(x,x+y)$ with a divisor in $(n,2n)$ is at most $\epsilon y$?




It is not clear what the intended quantifier on $x$ is. Cambie has observed that if this is intended to hold for all $x$ then, provided\[\epsilon(\log n)^\delta (\log\log n)^{3/2}\to \infty\]as $n\to \infty$, where $\delta=0.086\cdots$, there is no such $y$, which follows from an averaging argument and the work of Ford \cite{Fo08}. On the other hand, Cambie has observed that if $\epsilon\ll 1/n$ then $y(\epsilon,n)\sim 2n$: indeed, if $y&#60;2n$ then this is impossible taking $x+n$ to be a multiple of the lowest common multiple of $\{n+1,\ldots,2n-1\}$. On the other hand, for every fixed $\delta\in (0,1)$ and $n$ large every $2(1+\delta)n$ consecutive elements contains many elements which are a multiple of an element in $(n,2n)$.




References

[Fo08] Ford, Kevin, The distribution of integers with a divisor in a given interval . Ann. of Math. (2) (2008), 367-433.


\bigskip
\noindent\small{Source: \url{https://www.erdosproblems.com/450}}
\end{document}
